% AG.36: independent draft, integrated only after audit.
\subsection{A geometric unit-root obstruction along a fixed prime-power sequence}

\begin{claim}{AG.36}{Proved}\label{clm:AGfixedthreepower}
For each \(n\geq1\), choose a primitive \(3^n\)-th root
\(a_n\in\overline{\mathbb F}_2\), and put
\[
 d_n=2\cdot3^{n-1},\qquad k_n=\mathbb F_{2^{d_n}},\qquad
 E_n:\ y^2+xy=x^3+a_n.
 \tag{AG.36.1}
\]
Let \(\beta_n\in\mathbb Z_2^\times\) be the unit root of
\(X^2-t_nX+2^{d_n}\), where \(t_n\) is the Frobenius trace
of \(E_n/k_n\). Then
\[
 \lim_{n\to\infty}v_2(\beta_n-1)=+\infty,
 \qquad
 \exp E_n(k_n)[2^\infty]=2^{v_2(\beta_n-1)}.
 \tag{AG.36.2}
\]
Every \(E_n\) is a fibre of the single smooth ordinary elliptic
family \(y^2+xy=x^3+a\) over \(\mathbb G_{m,\mathbb F_2}\).
The parameter \(a_n\) has exact degree \(3^{n-1}\) over
\(\mathbb F_4\), an odd degree, and has exact multiplicative
order \(3^n\). Thus a fixed geometric family, weight-one Weil
purity, integrality, ordinariness, a fixed prime-power torsion
sequence, and bounded \(2\)-part of the parameter degrees do
not imply a uniform bound for the return valuations.

The conclusion asserts divergence, without a quantitative
comparison between \(2^{v_2(\beta_n-1)}\) and a prescribed
power of \(3^n\). It neither identifies this family with the
relative factors of an Abel tower nor disproves the weaker
polynomial growth bound required by that route.
\end{claim}

\begin{proof}
We first prove a finite-cover statement. If \(Y\to
\mathbb A^1_{\mathbb F_2}\) is a finite etale torsor under a
finite abelian \(2\)-group and its fibre at \(0\) is split,
then its fibre at every primitive \(3^n\)-th root is split over
\(k_n\) for all sufficiently large \(n\).
It is enough to prove this for a cyclic quotient of order
\(2^r\), and then take the maximum of finitely many bounds.
The Artin--Schreier--Witt exact sequence and affine vanishing
of \(H^1(\mathbb A^1,W_r\mathcal O)\) give
\[
 H^1_{\mathrm{\acute et}}(\mathbb A^1_{\mathbb F_2},
                         \mathbb Z/2^r)
 =W_r(\mathbb F_2[a])/(F-1)W_r(\mathbb F_2[a]).
 \tag{AG.36.3}
\]
Here \(F\) is Witt Frobenius. These standard exact statements
are recorded explicitly in \citet[Theorem~1.2 and the following
cohomology sequence, p.~2]{AG_Sekiguchi2001}; compare the
field-theoretic treatment and Frobenius trace formula in
\citet[Theorem~3.2, Section~5.4, p.~29]{AG_KostersWan2016}.
Affine vanishing also follows directly by induction from the
filtration of \(W_r\mathcal O\) with additive-group quotients.

Choose a representative \(w(a)\) in (AG.36.3).
On \(W_r(\mathbb F_2)\), Frobenius is the identity. Since
splitting at \(0\) says that the specialized cohomology class
vanishes, it follows that \(w(0)=0\), as an actual Witt vector.
There is a finite additive expansion
\[
                 w(a)=\sum_{i=0}^{r-1}\sum_{j\in J_i}
                                  V^i[a^j],
                 \qquad J_i\subset\mathbb Z_{>0}\text{ finite}.
 \tag{AG.36.4}
\]
Repeated terms may be retained in this expression. To obtain
it, subtract the sum of the Teichmuller monomials of the first
Witt coordinate. The difference has first coordinate zero,
so is uniquely \(V\) of a shorter Witt vector with polynomial
coordinates. This shorter vector still specializes to zero
at \(a=0\). Induction proves (AG.36.4), with no constant
monomial and only finitely many exponents. Let \(D\) be the
largest exponent occurring, taking \(D=0\) if the class is
zero.

The order of \(2\) modulo \(3^n\) is \(d_n\): the order
is even and \(v_3(4^m-1)=1+v_3(m)\). Hence the absolute
Frobenius conjugates of \(a_n\) are precisely all primitive
\(3^n\)-th roots. In the full Witt ring of \(k_n\), the
Teichmuller lift \(z=[a_n]\) has exact order \(3^n\), and
\[
 \sum_{s=0}^{d_n-1}z^{j2^s}
 =\sum_{\zeta\text{ primitive of order }3^n}\zeta^j=0
 \qquad(0<j<3^{n-1}).
 \tag{AG.36.5}
\]
Indeed the last sum is the sum over all \(3^n\)-th roots
minus the sum over all \(3^{n-1}\)-th roots; both geometric
sums vanish for the stated exponents. Since \(F\) commutes
with \(V\), (AG.36.4)--(AG.36.5) show that
\[
 \operatorname{Tr}_{W_r(k_n)/W_r(\mathbb F_2)}w(a_n)=0
                       \qquad\text{if }3^{n-1}>D.
 \tag{AG.36.6}
\]
For \(F(b)-b=w(a_n)\), the action of the degree-\(d_n\)
Frobenius on \(b\) is translation by this trace, by telescoping
\(F^{d_n}(b)-b\). Thus (AG.36.6) is exactly the splitting
condition. This proves the finite-cover statement.

We now apply it to the elliptic family. Its discriminant is
\(a\); each fibre over \(\mathbb G_m\) is ordinary and
has geometric etale \(2\)-rank one. For each \(r\geq1\),
let
\[
 Q_r=\ker\bigl(V^r:E^{(2^r)}\longrightarrow E\bigr).
 \tag{AG.36.7}
\]
This is a finite etale group scheme over \(\mathbb G_m\),
with geometric fibres cyclic of order \(2^r\). Its frame
torsor has group \((\mathbb Z/2^r)^\times\), a finite
abelian \(2\)-group. We show directly that this torsor is
trivial over the local field \(K=\mathbb F_2((a))\).
This statement concerns the etale group (AG.36.7), and does
not presume rationality of the non-etale group scheme
\(E[2^r]\) over the imperfect field \(K\).

In the purely inseparable extension \(K^{1/2}\), the point
\(R_1=(0,a^{1/2})\) has exact order two. Suppose a point
\(R=(x,y)\) of order \(2^s\) is defined over
\(K^{1/2^s}\), and both coordinates have positive valuation
(the zero coordinate is allowed). The convergent series
\[
                  \lambda=\sum_{i\geq0}x^{2^i}
 \tag{AG.36.8}
\]
belongs to the same complete local field, has positive
valuation, and satisfies \(\lambda^2+\lambda=x\).
Set
\[
 u^2=y+(\lambda+1)x,\qquad v=u(\lambda+u).
 \tag{AG.36.9}
\]
These coordinates belong to \(K^{1/2^{s+1}}\) and have
positive valuation. The curve equation for \(R\) gives
\(u^4+xu^2=a\); in particular \(u\ne0\).
It follows directly that \(v^2+uv=u^3+a\), and the
elliptic doubling formulas give \(2(u,v)=R\).
Thus \((u,v)\) has exact order \(2^{s+1}\).
Induction constructs a generator of the geometric
\(2^r\)-torsion in \(E(K^{1/2^r})\).
Applying relative Frobenius \(F^r\) gives a \(K\)-rational
generator of \(Q_r\). Consequently \(Q_r/K\) is constant.

The frame torsor is therefore unramified at \(a=0\),
with trivial residue Frobenius there. Its normalization over
\(\mathbb A^1_{\mathbb F_2}\) is finite etale and has split
fibre at \(0\). Apply the finite-cover statement to this
extension. For each fixed \(r\), and all sufficiently large
\(n\), the group \(Q_r\) at \(a_n\) is constant over
\(k_n\). Relative Frobenius identifies its geometric points
Galois-equivariantly with those of \(E_n[2^r]\), so the
unit root satisfies \(\beta_n\equiv1\pmod{2^r}\).
This proves the limit in (AG.36.2).

Finally, the rational \(2\)-primary subgroup of an ordinary
elliptic curve over a finite field is cyclic. The other
Frobenius root \(2^{d_n}/\beta_n\) is divisible by two,
and therefore
\[
 v_2\#E_n(k_n)
 =v_2\bigl((1-\beta_n)
              (1-2^{d_n}/\beta_n)\bigr)
 =v_2(1-\beta_n).
\]
This proves the exponent formula. The exact parameter degree
over \(\mathbb F_4\) is \(\operatorname{ord}_{3^n}(4)
=3^{n-1}\), as asserted.
\end{proof}

\paragraph{Finite certificates and quantitative scope.}
The independent complete point-count program
\texttt{agent\_geometric/code/kloosterman\_three\_power\_test.cpp}
and its JSON output use the irreducible polynomial
\(\Phi_{3^n}(X)=X^{2\cdot3^{n-1}}+X^{3^{n-1}}+1\)
for \(n=1,2,3\). They verify, respectively,
\[
\begin{array}{c|r|r|r|r}
 n&\#k_n&\#E_n(k_n)&t_n&v_2(\beta_n-1)\\\hline
 1&4&4&1&2\\
 2&64&64&1&6\\
 3&262144&262912&-767&8
\end{array}
\]
The trace-one pattern of the first two examples already
fails at the third. The proof of divergence does not use
these finite data. Its threshold for a fixed torsion level
is the finite degree bound \(3^{n-1}>D\) in (AG.36.6),
where \(D\) comes from representatives of that level's
frame-torsor characters. No bound for the growth of these
representative degrees with the torsion level is asserted.
